\documentclass[12pt]{article}
\usepackage{ctex}
\usepackage{amsmath, amssymb}
\usepackage{geometry}
\geometry{a4paper, margin=2.5cm}
\usepackage{enumitem}
\title{第 6 章\ 刚体力学\ 例题}
\author{}
\date{}

\begin{document}
\maketitle

\begin{enumerate}[label=\textbf{例\arabic*.}, leftmargin=*]

% 例 1
\item 一质量为 $m$、长为 $L$ 的均匀细杆，初始静止于水平位置，可绕通过其端点 $O$ 且垂直于杆的固定轴在竖直平面内无摩擦地转动。求：(1) 杆在水平位置时开始释放瞬间的角加速度；(2) 杆转到竖直位置时的角速度。

% 例 2
\item 已知棒长 $L$、质量 $M$，放在摩擦系数为 $\mu$ 的水平桌面上，绕其端点 $O$ 转动（如图）。求摩擦力对 $O$ 点的力矩。

% 例 3
\item 半径为 $R$、长为 $L$、质量为 $M$ 的实心圆柱体，求对中心直径（过质心且垂直于柱轴）的转动惯量。

% 例 4
\item 半径为 $R$、质量为 $M$ 的匀质圆盘，求绕过盘边缘且与盘面中心轴平行的轴的转动惯量。

% 例 5
\item 在半径为 $R$、质量为 $M$ 的匀质圆盘上，挖出半径为 $r$ 的两个圆孔，圆孔中心在圆盘半径的中点。求剩余部分对大圆盘中心且与盘面垂直的轴线的转动惯量。

% 例 6
\item 滑轮半径为 $R$，质量为 $M$（视为匀质圆盘），绳子不可伸缩的轻绳，绳子与滑轮间无滑动，轴处无摩擦，两个悬挂物质量分别为 $m_1$、$m_2$。求：两重物的加速度、滑轮的角加速度、绳中的张力。

% 例 7
\item 两个皮带轮半径分别为 $R_1$、$R_2$，质量分别为 $m_1$、$m_2$，分别绕固定轴 $O_1$、$O_2$ 转动，用皮带相连，轮 1 作用力矩 $M_1$，轮 2 有负载力矩 $M_2$，皮带与轮无滑动，轴处无摩擦。求轮 1 的角加速度。

% 例 8
\item 长为 $l$、质量分布不均匀的细杆，其线密度为 $\lambda = a + br$（$a$、$b$ 为常量），该杆可绕 $A$ 端的 $z$ 轴在铅锤面内转动。现将杆从水平位置（$\theta=\pi/2$）释放，求杆转到铅锤位置（$\theta=0$）过程中重力做的功。

% 例 9
\item 圆盘滑轮质量 $M$、半径 $R$，绕轻绳，绳的另一端系一质量 $m$ 的物体，轴无摩擦，开始时系统静止。求物体下降 $s$ 时滑轮的角速度 $\omega$ 与角加速度 $\beta$。

% 例 10
\item 以力 $F$ 将一块粗糙平面均匀压在轮上，平面与轮之间的滑动摩擦系数为 $\mu$，轮为匀质圆盘，半径为 $R$，质量为 $M$，轴处摩擦力不计，轮的初角速度为 $\omega_0$。问：轮转过多少度时即停止转动。

% 例 11
\item 半径为 $R$、质量为 $4m$ 且均匀分布的转盘可绕铅垂轴无摩擦地转动，盘上站有质量为 $m$ 的四个人：两个人站在盘的边沿，两个人站在 $R/2$ 处，整个系统开始以角速度 $\omega_0$ 转动。此后，人相对于盘作圆周运动，站在边沿上的两人相对于盘的速度为 $u$，站在 $R/2$ 处的两人相对于盘的速度为 $2u$。求：(1) 人走动后转盘的角速度；(2) 要使转盘停止转动，$u$ 的大小。

% 例 12
\item 长为 $L$、质量为 $M$ 的均匀杆，一端悬挂，由水平位置无初速度地下落，在铅直位置与质量为 $m$ 的物体 $A$ 做完全非弹性碰撞，碰后物体 $A$ 沿摩擦系数为 $\mu$ 的水平面滑动。求物体 $A$ 滑动的距离。

% 例 13
\item 半径为 $R$ 的转盘可绕过中心 $O$ 的铅垂轴无摩擦转动。一子弹质量 $m$、速度 $v$ 沿水平方向射入静止的转盘并嵌入盘边。设转盘质量 $M$、半径 $R$。求：子弹嵌入后转盘的角速度。

% 例 14
\item 一质量为 $m$、长为 $l$ 的均质细杆，可绕通过中心的固定水平轴在铅垂面内自由转动。开始时杆静止于水平位置。一质量与杆相同的昆虫以速度 $v_0$ 垂直落到距点 $O$ 为 $l/4$ 处的杆上，昆虫落下后立即向杆的端点爬行。若要使杆以匀角速度转动，求昆虫沿杆爬行的速度。

% 例 15
\item 一长为 $L=1\,\mathrm{m}$ 的均质细杆，可绕通过一端的水平光滑轴 $O$ 在铅垂面内自由转动。开始时杆处于铅垂位置，今有一粒子弹沿水平方向以 $v=10\,\mathrm{m/s}$ 的速度射入细杆。设入射点离 $O$ 点的距离为 $\frac{3L}{4}$，子弹的质量为杆质量的 $\frac{1}{9}$。求：(1) 子弹与杆开始共同运动的角速度；(2) 子弹与杆共同摆动能达到的最大角度。

\end{enumerate}

\end{document}
